\chapter{Auditoría del desarrollo}
\label{anexo:desarrollo}

Este anexo conserva solo los defectos que pudieron cambiar la validez o la interpretación del
estudio. La historia completa permanece en la bitácora del repositorio; repetirla aquí añadiría
cronología, no evidencia. La Tabla~\ref{tab:defectos} resume síntoma, riesgo y corrección.

\begin{table}[H]
  \centering
  \footnotesize
  \caption{Defectos del desarrollo con impacto sobre la validez.}
  \label{tab:defectos}
  \resizebox{\textwidth}{!}{% Fuente: docs/bitacora.md, entrada 2026-07-28 «Corrección de validez, alfa neto y limpieza».
% Tabla transcrita del registro de desarrollo, no generada por script: su origen es la
% bitácora del proyecto y no un artefacto del Study.
\begin{tabularx}{\textwidth}{XX}
\toprule
Defecto detectado & Efecto medido de la corrección \\
\midrule
La puerta de no inferioridad penalizaba a los candidatos \textbf{superiores}: cuanto mayor la
diferencia, más ancho el intervalo &
\texttt{feature\_preset} pasa de \texttt{core} (Rank-IC 0,0730) a \texttt{all} (0,0958), con
ventaja pareada $+0{,}0216$. Ningún retador era elegible pese a dominar \\
\addlinespace
\texttt{market\_regime\_feature} se decidió sobre una ventaja de 0,00112, por debajo del ruido &
Pasa de \texttt{True} a \texttt{False} por simplicidad \\
\addlinespace
\texttt{meta\_method} se decidió sobre una ventaja de 0,00033 &
Se mantiene, pero registrado como empate técnico y no como victoria \\
\addlinespace
Emparejamiento por cadena de fecha con rejillas desplazadas por \texttt{execution\_lag\_days} &
Se empareja por periodo mensual; sin bloque completo común el resultado se marca no aplicable en
vez de devolver \texttt{ci\_low = 0,0} \\
\addlinespace
\texttt{evaluation\_key} no incluía el hash del dataset &
La misma clave aparecía con CAGR 0,1468 y 0,1692. Ahora la clave separa datasets \\
\addlinespace
Seis factores de precio se inyectaban en el agente \textit{momentum} fuera de todo condicional &
La ablación \texttt{fundamental} deja de recibir información de precio y mide lo que declara \\
\addlinespace
Dos definiciones incompatibles de \textit{information ratio} bajo el mismo nombre &
Una sola, anualizada \\
\addlinespace
CAGR, \textit{drawdown} y alfa mezclaban la era reservada con la de selección &
Métricas segmentadas: selección, confirmación y curva completa \\
\addlinespace
\texttt{geometric\_excess\_return} era una resta de CAGR &
Cociente de acumulados \\
\addlinespace
Una posición sin precio se marcaba plana &
Convención de exclusión tipo CRSP ($-30\,\%$) y liquidación contra efectivo \\
\addlinespace
\texttt{subsample=0.8} sin \texttt{subsample\_freq}: el \textit{bagging} nunca se activaba &
\texttt{subsample\_freq=1} \\
\addlinespace
Nulo de carteras aleatorias con percentil 95 de CAGR del 107\,\% &
Exige cobertura anual completa, aplica la guarda de datos y paga los mismos costes \\
\bottomrule
\end{tabularx}
}
\end{table}

\section*{Tres lecciones reutilizables}

\textbf{1. Declarar no equivale a calcular.} Varias opciones del catálogo llegaron a existir en la
interfaz antes de modificar realmente las variables consumidas por el modelo. El sistema podía
terminar sin error y producir métricas plausibles. La corrección fue convertir la correspondencia
entre configuración declarada, datos preparados y evaluación en un contrato automático.

\textbf{2. Un resultado plausible puede esconder un fallo semántico.} Una puerta de no inferioridad
penalizaba candidatos superiores por tener intervalos más anchos; una calibración degenerada podía
emitir expectativas constantes; y una posición sin puntuación detenía una rotación completa. Ninguno
de esos casos necesitaba lanzar una excepción. Los diagnósticos deben comprobar invariantes y no solo
que el proceso finalice.

\textbf{3. Toda decisión económica necesita contrafactual y contabilidad.} Una venta solo es
interpretable si existe un destino mejor después de costes, y una comisión solo puede cobrarse si la
orden se ejecuta. Esta regla obligó a unir señal, orden, posición y coste en el mismo registro y
evitó que un backtest aparentemente conservador destruyera valor mediante operaciones inexistentes.

La consecuencia para el TFM es directa: los artefactos finales no se aceptan por parecer razonables,
sino porque los contratos permiten reconstruir qué dato estaba disponible, qué decisión se tomó y
qué efecto contable produjo. Las correcciones anteriores explican la necesidad de esa disciplina;
no se presentan como evidencia económica adicional.
